\documentclass[11pt,a4paper]{article}
\usepackage[T1]{fontenc}
\usepackage[utf8]{inputenc}
\usepackage{lmodern}
\usepackage[margin=27mm]{geometry}
\usepackage{amsmath,amssymb,amsthm,mathtools,booktabs,array}
\usepackage{microtype}
\usepackage[hidelinks]{hyperref}
\usepackage{xurl}
\hypersetup{pdflang={en-US}}
\hypersetup{pdftitle={Covariance Signs for Residue-Class Occupancy on Decimal Unit Carriers},pdfauthor={Oleksiy Babanskyy}}
\newtheorem{theorem}{Theorem}
\newtheorem{lemma}[theorem]{Lemma}
\newtheorem{proposition}[theorem]{Proposition}
\newtheorem{corollary}[theorem]{Corollary}
\theoremstyle{remark}
\newtheorem{remark}[theorem]{Remark}
\DeclareMathOperator{\Cov}{Cov}
\DeclareMathOperator{\sgn}{sgn}
\newcommand{\E}{\mathbb E}
\newcommand{\one}{\mathbf 1}
\newcommand{\ff}[2]{(#1)_{#2}}
\newcommand{\vect}[3]{(#1,#2,#3)}
\numberwithin{equation}{section}
\setlength{\emergencystretch}{2em}
\title{Covariance Signs for Residue-Class Occupancy\\on Decimal Unit Carriers}
\author{Oleksiy Babanskyy}
\date{}
\begin{document}
\maketitle
\begin{abstract}
Choose a uniformly random fixed-size subset of the integers below a power
of ten that are coprime to ten. We determine the sign of the covariance
between the numbers of occupied residue classes modulo \(3\) and modulo
any positive integer coprime to \(30\). Outside deterministic zero ranges,
the covariance is strictly negative except for four families of moduli
near half the cutoff. Explicit quadratic expressions determine every
exceptional sign; a divisibility argument excludes nontrivial zeros in
two of the four families. We also prove a sufficient
negative-covariance criterion for arbitrary partitions with three nearly
equal rows and controlled column profiles, with a consequence for
nonnegative weighted sums over overlapping admissible test sets.
The proof combines empty-class
identities, a periodic matching count, majorization, and exact polynomial
certificates checked by rational interpolation. A minimal standard-library
companion verifies all certificate identities. For prime moduli, the
classification gives the sign of the covariance of two local sieve
factors. It does not assert a distribution law for actual prime sets.
\end{abstract}

\section{Introduction and main results}\label{sec:intro}

The integers below \(100\) that are coprime to \(10\) form the familiar
four-column display
\[
 \{10j+r:0\le j\le9,\ r\in\{1,3,7,9\}\}.
\]
Every prime at least \(11\) has one of these four terminal digits.
The display is therefore a natural carrier for studying local divisibility
constraints. Its entries, however, include composite integers and the
integer \(1\). The carrier specifies candidates; it does not specify their
primality. This distinction leads to a finite question that can be answered
without a prime-distribution hypothesis: when a fixed number of candidate
offsets are sampled, how do two counts of occupied residue classes vary
together?

Let \(d\ge2\), \(N=10^d\), and
\begin{equation}\label{eq:carrier}
 \Omega_N=\{h\in\mathbb Z:1\le h<N,\ \gcd(h,10)=1\},\qquad
 m=|\Omega_N|=\frac{2N}{5}=3a+4,\qquad s=2a+4.
\end{equation}
Here
\[
 a=\frac{4(10^{d-1}-1)}3,
 \qquad a=12\ \text{if }d=2,\quad a\ge132\ \text{if }d\ge3.
\]
For a uniformly chosen \(k\)-element subset \(H\) of \(\Omega_N\), set
\[
 \nu_q(H)=|\{h\bmod q:h\in H\}|,\qquad
 C(d,q,k)=\Cov(\nu_3(H),\nu_q(H)).
\]
Unless stated otherwise, \(0\le k\le m\), \(q\) is a positive integer,
and \(\gcd(q,30)=1\). The modulus \(q\) need not be prime.
Binomial coefficients with nonnegative upper index are zero when the
lower index is outside its usual range.

\begin{theorem}[Decimal sign classification]\label{thm:decimal}
If \(q=1\) or \(q\ge N/2\), then \(C(d,q,k)=0\) for every \(k\).
For \(7\le q<N/2\), the covariance is zero when \(k=0,1\) or \(k>s\).
For \(2\le k\le s\), it is strictly negative except possibly when
\[
 q=N/2-g,\qquad
 g\in\{1,3,7\},\quad\text{or}\quad g=19\ \text{and }d\ge3.
\]
In these exceptional families put \(X=(s-k)(s-k-1)\). The sign of \(C\),
including equality, is the sign of the following expression:
\begin{center}
\begin{tabular}{@{}cl@{}}
\toprule
Gap \(g\)&Sign expression\\
\midrule
\(1,3\)&\(a(5a+7)-2X\)\\
\(7\)&\((2a+1)(7a^2-7a-24)-(11a-2)X\)\\
\(19,\ d\ge3\)&\((2a+1)(a^2-49a-72)-(17a-14)X\)\\
\bottomrule
\end{tabular}
\end{center}
Each exceptional family starts with negative covariance at \(k=2\) and
ends with positive covariance at \(k=s-1,s\). It has one transition
from negative to positive and at most one intervening integer zero.
\end{theorem}

The statement is an exact sign decision for every admitted input.
Corollary~\ref{cor:nonzero} rules out nontrivial integer zeros for gaps
\(7\) and \(19\). The corresponding arithmetic question remains open
only for gaps \(1\) and \(3\); it is separate from deciding the sign
at a specified input.

Our second result has no arithmetic carrier in its hypotheses.
For a partition, its occupancy is the number of its classes met by the
random subset. A column profile records its intersections with the three
rows, in the specified row order.

\begin{theorem}[A three-row criterion]\label{thm:criterion}
Let a finite set have a row partition of sizes
\(\rho=(a+2,a+2,a)\), where \(a\ge40\) is an integer, and a partition
into nonempty columns. For each column write
\(v=b\one+e\), where \(\one=(1,1,1)\), \(b=\min_i v_i\), and \(\min_i e_i=0\).
Suppose every column satisfies
\begin{equation}\label{eq:criterion}
 b\ge4,\qquad \max_i e_i\le2,\qquad |v|\le a,\qquad
 e=(0,0,2)\ \Longrightarrow\ b\ge6.
\end{equation}
For a uniform \(k\)-subset, the covariance of the two occupancies is
strictly negative for \(2\le k\le2a+4\), and zero for the other
cardinalities \(0\le k\le3a+4\).
\end{theorem}

The orientation of \((0,0,2)\) is essential: the third row is the smaller
one. The criterion is sufficient, with no claim that its numerical
thresholds are optimal. It is nonvacuous for infinitely many populations.
For example, take \(a=12t\), \(t\ge4\), with \(3t-1\) columns of profile
\((4,4,4)\) and one column of profile \((6,6,4)\).
Their sum is \((a+2,a+2,a)\) and all hypotheses hold.
The argument controls each admissible set separately.
Corollary~\ref{cor:weighted} therefore permits overlapping test sets
and fixed nonnegative weights, without changing the sampling law.

\subsection{Relation to established frameworks}\label{sec:context}
The elementary empty-class formula used below is a finite
hypergeometric identity. The occupancy of a partition is also the rank
function of the direct sum of rank-one uniform matroids on its classes.
Consequently the joint generating function
\[
 \sum_{H\subseteq\Omega_N} z^{|H|}u^{\nu_3(H)}v^{\nu_q(H)}
\]
is an instance of the established rank generating function of a pair of
matroids; see Las Vergnas \cite[p.~305]{lasvergnas}. No new general
matroid invariant is proposed here. The contribution is the sign
classification for this structured pair, and the separate sufficient
criterion in Theorem~\ref{thm:criterion}.

Negative association is an important framework for sampling without
replacement \cite{joagdev}. Its disjoint-coordinate condition does not
directly decide our covariance: both occupancies depend on the same
sample indicators, and a row event and a column event generally overlap.
Indeed Theorem~\ref{thm:decimal} contains positive covariances.
Scalar hypergeometric stochastic-order criteria \cite{klm} and
coupling comparisons for weighted sums \cite{benhamou} concern related
sampling questions. The urn comparison needed here changes three
population sizes simultaneously and adds three empty-row probabilities;
we prove the required comparison explicitly rather than assuming a
componentwise stochastic order.

The arithmetic motivation is equally specific. Residue occupancies
determine the local factors occurring in singular-series calculations;
constrained averages of singular series are studied, for example, by
Kuperberg \cite{kuperberg}. Section~\ref{sec:sieve} gives the exact
two-factor consequence of Theorem~\ref{thm:decimal}. An average of an
entire singular series, or a statement about actual prime patterns,
requires additional information.

\subsection{Proof organization and computational content}
We first reduce occupancy covariance to empty-class probabilities.
Two certified urn inequalities and majorization prove
Theorem~\ref{thm:criterion}, which controls the small-modulus range.
The remaining arithmetic problem splits into a matching window and
fourteen finite bands. The matching window gives the exceptional
families explicitly. In each band, affine column multiplicities reduce
the sign to finitely many polynomial identities.

All load-bearing certificate data are supplied with the article.
The verifier uses only integer arithmetic and exact rational arithmetic
from Python's standard library. It reconstructs the polynomials and
checks identities on complete interpolation grids whose degree bounds
are proved below. These checks establish polynomial identities on
infinite parameter domains; they are not numerical extrapolations.
The certificate specification appears in Section~\ref{sec:certificates}.
The proof is computer-assisted, not a proof-assistant formalization.

\section{Empty-class identities}\label{sec:empty}
A complete interval of length \(30\) contributes four carrier points
to each residue class modulo \(3\). Since \(N\equiv10\pmod{30}\),
the final interval contributes \((2,2,0)\). Thus the row sizes are
\begin{equation}\label{eq:rows}
 \rho=(a+2,a+2,a)
\end{equation}
in the fixed order of residues \(0,1,2\).

\begin{lemma}[Covariance from intersections]\label{lem:empty}
For two partitions of a set of size \(m\), let \(\rho_r\) be the row
sizes, \(j_c\) the column sizes and \(v_{rc}\) the intersection sizes.
For uniform \(k\)-subsets, with \(Z=\binom mk\), their occupancy
covariance equals
\begin{equation}\label{eq:empty}
 \sum_{r,c}\left\{
 \frac{\binom{m-\rho_r-j_c+v_{rc}}k}{Z}
 -\frac{\binom{m-\rho_r}k\binom{m-j_c}k}{Z^2}\right\}.
\end{equation}
\end{lemma}
\begin{proof}
Each occupancy is the total number of classes minus the sum of the
indicators that a class is empty. The covariance is therefore the sum
of the covariances of empty-row and empty-column indicators.
Avoiding their union leaves \(m-\rho_r-j_c+v_{rc}\) available elements,
which proves the formula.
\end{proof}

For a triple \(R=(R_0,R_1,R_2)\) of nonnegative population sizes, put
\begin{equation}\label{eq:F}
 F(R;k)=\sum_{r=0}^2\frac{\binom{M-R_r}k}{\binom Mk},
 \qquad M=R_0+R_1+R_2,\quad 0\le k\le M.
\end{equation}
This is the expected number of empty rows.
For a column \(v\) of size \(j=|v|\), conditioning on its being empty
gives its contribution to \eqref{eq:empty} as
\begin{equation}\label{eq:conditional}
 C_v=\frac{\binom{m-j}k}{\binom mk}
       \bigl(F(\rho-v;k)-F(\rho;k)\bigr),\qquad k\le m-j.
\end{equation}
For \(k>m-j\) that column is always occupied, so its contribution is zero.
Column multiplicities must be retained when summing \(C_v\).

The deterministic boundaries are immediate. When \(k=0,1\), both
occupancies are deterministic. If \(k>s=m-a\), every row is occupied.
When \(q=1\), the column occupancy is deterministic. Finally, if
\(q\ge N/2\), two distinct carrier points cannot be congruent modulo \(q\).
Both points are odd and \(q\) is odd, so their nonzero difference would
have magnitude at least \(2q\ge N\). Hence \(\nu_q(H)=k\).

\subsection{A common polynomial kernel}\label{sec:kernel}
Write \((z)_j=z(z-1)\cdots(z-j+1)\), with \((z)_0=1\).
For a collection of column profiles choose
\[
 J\ge\max_v|v|,\qquad
 D\ge\max_{v,r}\delta_r(v),\qquad
 \delta_r(v)=|v|-v_r+2\one_{\{r<2\}},
\]
and assume \(D\ge2\). Define
\begin{align}
 K_v(a,k)={}&(m)_J\sum_{r=0}^2(s-k)_{\delta_r(v)}
                     \prod_{i=\delta_r(v)}^{D-1}(s-i) \notag\\
 &-\left(\prod_{i=2}^{D-1}(s-i)\right)
       \bigl[s(s-1)+2(s-k)_2\bigr](m-k)_{|v|}
                     \prod_{i=|v|}^{J-1}(m-i).
 \label{eq:kernel}
\end{align}
Empty products are one. Whenever \(2\le k\le s\), \(J\le m\),
and \(D\le s\), the empty-class identity yields
\begin{equation}\label{eq:kernelcov}
 C=\frac{\binom sk}{\binom mk(s)_D(m)_J}\sum_v n_v K_v(a,k),
\end{equation}
where \(n_v\) is the number of columns with profile \(v\).
Indeed \(m-\rho_r-|v|+v_r=s-\delta_r(v)\), and
\[
 \frac{\binom{s-\delta}k}{\binom sk}
       =\frac{(s-k)_\delta}{(s)_\delta},\qquad
 \frac{\binom{m-j}k}{\binom mk}=\frac{(m-k)_j}{(m)_j}.
\]
Moreover
\(F(\rho;k)=\binom sk\bigl[s(s-1)+2(s-k)_2\bigr]/
 \bigl(\binom mk s(s-1)\bigr)\).
These identities also hold at integer values where the numerator
binomial coefficient vanishes.

The raw total degree of \(K_v\) in \((a,k)\) is at most \(J+D\),
and its degree in \(k\) is at most \(J+2\). Each first summand has
total degree \(J+D\) and \(k\)-degree \(\delta_r(v)\le |v|+2\le J+2\);
the second summand has the same total bound and \(k\)-degree
at most \(|v|+2\).

\section{Two finite certificates for urn comparisons}\label{sec:urn}
We state the certificate principle explicitly, so that the role of
each finite computation is clear.

\begin{lemma}[Identity and sign certificates]\label{lem:certificate}
Two polynomials of degree at most \(A\) in one variable and at most
\(B\) in another are identical if they agree on a product of \(A+1\)
distinct nodes and \(B+1\) distinct nodes.
If
\[
 Q(t,x)=\sum_{i=0}^n B_i(t)\binom ni x^i(1-x)^{n-i}
\]
and each \(L B_i(t_0+y)\), for a positive integer \(L\), has
nonnegative power coefficients and a strictly positive constant term,
then \(Q(t,x)>0\) for \(t\ge t_0\), \(0\le x\le1\).
The analogous statement with \(-L B_i\) gives strict negativity.
\end{lemma}
\begin{proof}
Apply the univariate root bound successively in the two variables.
For the sign statement, every coefficient polynomial is strictly
positive for \(y\ge0\), and the Bernstein basis is nonnegative
and has sum one on \([0,1]\).
\end{proof}

\begin{lemma}[Balanced urn increment]\label{lem:increment}
Let \(R=(h,h+p,h+r)\), where \(h\ge16\) and \(0\le p\le r\le4\)
are integers. In the common domain \(0\le k\le3h+p+r\),
\[
 F(R+\one;k)\ge F(R;k).
\]
The inequality is strict for \(2\le k\le2h+p+r+2\), and is an
equality for the other cardinalities in that domain.
\end{lemma}
\begin{proof}
Put \(M=3h+p+r\), \(n=2h+p+r\), \(D=r+2\).
The list \(\mathcal D=(0,p,r)\) retains repetitions. Set
\begin{align*}
 A&=\sum_{d\in\mathcal D}(n+2-k)_d
                      \prod_{i=d}^{D-1}(n+2-i),\\
 B&=\sum_{d\in\mathcal D}(n+2-k)_{d+2}
                      \prod_{i=d+2}^{D-1}(n+2-i),\\
 P&=(M+3-k)_3 A-(M+3)_3 B.
\end{align*}
Conversion of the binomial coefficients in \eqref{eq:F} gives
\begin{equation}\label{eq:increment}
 F(R+\one;k)-F(R;k)
 =\frac{\binom{n+2}k}{\binom{M+3}k}
       \frac{P}{(n+2)_D(M+3-k)_3}
\end{equation}
for \(2\le k\le n+2\). Every factor outside \(P\) is positive
in this range: \(n+2\ge D\) and \(M+3-k\ge h+1\ge17\).

For each of the fifteen pairs \((p,r)\), the certificate
\texttt{urn\_monotonicity.json} establishes
\[
 P=k(k-1)Q,\qquad
 Q\bigm|_{k=2+(2h+p+r)x}
   =\sum_{i=0}^{r+1} B_i(h)\binom{r+1}{i}x^i(1-x)^{r+1-i}.
\]
It records the coefficients of \(L B_i(16+y)\) for a positive
integer scale \(L\). They are nonnegative and their constants are
positive. The parameter degree bound of the stored coefficient
polynomials is \(r+3\); there are \(490\) stored integers in total.

For completeness, the raw degree bounds follow from the displayed
formula for \(P\): its total degree in \((h,k)\) is at most
\(r+5\), and its \(k\)-degree is at most \(r+3\).
After \(k=2+(2h+p+r)x\), its degree in \(h\) remains at most \(r+5\),
and its degree in \(x\) at most \(r+3\). The encoded right side
has these same bounds. The checker reconstructs both sides and
compares them at \(r+6\) values of \(h\), starting at \(16\),
and \(r+4\) values \(x=j/(r+3)\). All \(890\) equalities hold.
Lemma~\ref{lem:certificate} therefore proves the identities globally
and \(Q>0\) on the required domain. Equation~\eqref{eq:increment}
proves strictness.

For \(k>n+2\), both empty-row expectations vanish.
For \(k=0,1\), they are respectively \(3,2\), proving the remaining
equalities.
\end{proof}

\begin{lemma}[Six column anchors]\label{lem:anchors}
Let \(\rho=(a+2,a+2,a)\), \(a\ge40\). For \(2\le k\le2a+4\),
\[
 F(\rho-v;k)<F(\rho;k)
\]
for each of the ordered vectors
\begin{equation}\label{eq:anchors}
 v\in\{(4,4,4),(4,4,5),(4,5,5),(4,5,6),(4,6,6),(6,6,8)\}.
\end{equation}
\end{lemma}
\begin{proof}
For each anchor use \eqref{eq:kernel} with \(J=|v|\) and
\(D=\max_r\delta_r(v)\). Then
\begin{equation}\label{eq:anchoridentity}
 \frac{\binom{m-J}k}{\binom mk}
       \bigl(F(\rho-v;k)-F(\rho;k)\bigr)
 =\frac{\binom sk}{\binom mk}\frac{K_v(a,k)}{(s)_D(m)_J}.
\end{equation}
The certificate \texttt{column\_anchors.json} establishes
\(K_v=k(k-1)Q_v\). After \(k=2+(2a+2)x\), it gives the degree-\(J\)
Bernstein expansion of \(Q_v\), whose coefficient polynomials have
degree at most \(J+D-2\) in \(a\).
For each coefficient \(B_i(a)\), the power coefficients of
\(-L B_i(40+y)\) are nonnegative with positive constant term.
Thus the \emph{Bernstein coefficients themselves} are strictly negative
for \(a\ge40\).

There are \(2630\) stored integers. The raw bounds \(J+D\) in
total degree and \(J+2\) in \(k\) follow from
Section~\ref{sec:kernel}. They give bounds \(J+D,J+2\) after
the substitution. The encoded right side has the same bounds.
The checker verifies the raw identity on
\((J+D+1)(J+3)\) rational grid points for each anchor,
totalling \(3166\) exact equalities. Lemma~\ref{lem:certificate}
therefore proves the global identity and strict negativity.
All denominators in \eqref{eq:anchoridentity} are positive.
Also \(J\le20\le a\), so \(m-J\ge s\) and its left prefactor
is positive throughout the asserted range.
\end{proof}

\section{Proof of the three-row criterion}\label{sec:criterionproof}
At fixed total size \(M\), the function \(F(R;k)\) is Schur-convex
on integer triples: making the row sizes more unequal increases it
weakly. This follows directly from discrete convexity, since
\begin{equation}\label{eq:convex}
 \binom{M-z-2}k-2\binom{M-z-1}k+\binom{M-z}k
       =\binom{M-z-2}{k-2}\ge0
\end{equation}
whenever the indicated row-size comparison is defined.
Equivalently, a unit transfer from a smaller entry to a larger one
cannot decrease the sum in \eqref{eq:F}.

\begin{proof}[Proof of Theorem~\ref{thm:criterion}]
There are nineteen possible normalized profiles
\(e\in\{0,1,2\}^3\) with \(\min e=0\).
For each choose a comparison profile \(e^*\) and an anchor base
\(b^*\) from Table~\ref{tab:major}.

\begin{table}[htbp]
\centering
\begin{tabular}{@{}lccc@{}}
\toprule
Profile class&\(e^*\)&\(b^*\)&\(b^*\one+e^*\)\\
\midrule
\(|e|=0\)&\((0,0,0)\)&4&\((4,4,4)\)\\
\(|e|=1\)&\((0,0,1)\)&4&\((4,4,5)\)\\
\(|e|=2,\ e\ne(0,0,2)\)&\((0,1,1)\)&4&\((4,5,5)\)\\
\(|e|=3\)&\((0,1,2)\)&4&\((4,5,6)\)\\
\(|e|=4\)&\((0,2,2)\)&4&\((4,6,6)\)\\
\(e=(0,0,2)\)&\((0,0,2)\)&6&\((6,6,8)\)\\
\bottomrule
\end{tabular}
\caption{Comparison profiles for the three-row criterion.}\label{tab:major}
\end{table}

The vector \((2-e^*_0,2-e^*_1,-e^*_2)\) majorizes
\((2-e_0,2-e_1,-e_2)\): their sums agree, the former minimum is
no larger, and its maximum is no smaller. These are precisely the
two majorization inequalities for triples of equal sum.
The table can also be verified over the nineteen profiles by the
companion. Adding the common value \(a-b\) gives
\[
 F(\rho-b\one-e;k)\le F(\rho-b\one-e^*;k).
\]

Since \(3b\le |v|\le a\), the minimum of either residual population
is at least
\[
 a-b-2\ge 2a/3-2>16,
\]
and its spread is at most four. By symmetry of \(F\), the entries
may be sorted before applying Lemma~\ref{lem:increment}.
Decreasing \(b\) one unit at a time until \(b=b^*\) adds \(\one\)
at each step. The hypotheses give \(b\ge b^*\), so
\[
 F(\rho-v;k)
 \le F(\rho-b\one-e^*;k)
 \le F(\rho-b^*\one-e^*;k)
 < F(\rho;k).
\]
The final strict inequality is Lemma~\ref{lem:anchors}.
The fixed \(k\) lies in the domain of every intermediate urn,
because the initial remaining population has size \(m-|v|\ge s\).
Where necessary the increment lemma is used weakly, including ranges
where both empty-row expectations vanish.

For \(2\le k\le s\), the prefactor in \eqref{eq:conditional} is
positive for every column, again because \(m-|v|\ge s\).
Every contribution is strictly negative, so their sum is strictly
negative. The zero ranges follow from Section~\ref{sec:empty}.
\end{proof}

\begin{remark}[Why the orientation restriction is retained]\label{rem:obstruction}
For \(a=40\), \(k=2\), the column \(v=(4,4,6)\) has a
\emph{positive} contribution \(C_v=71/4846323\).
This can be checked by counting pairs: put
\(I=\mathbf1_{H\cap B\ne\varnothing}\) and \(D=\binom{124}{2}=7626\).
Across these pairs, \(\sum U=12750\), \(\sum I=1631\) and
\(\sum UI=2727\); hence
\[
 C_v=\frac{2727}{7626}
     -\frac{12750\cdot1631}{7626^2}
     =\frac{71}{4846323}>0.
\]
Thus the condition on the normalized profile \((0,0,2)\) cannot be
removed from this columnwise proof. This is a counterexample to
the individual-column inequality, not a counterexample for an
entire partition satisfying weakened hypotheses.
Similarly a completely unrestricted balanced urn increment is false:
\[
 F((2,6,6);2)-F((1,5,5);2)=-\frac{23}{1001}.
\]
\end{remark}

\section{Small moduli from periodic column profiles}\label{sec:small}
The criterion has an application slightly wider than the decimal carrier.

\begin{proposition}\label{prop:small}
Let \(N\equiv10\pmod{30}\), \(q\ge7\), \(\gcd(q,30)=1\), and \(N>32q\).
Use the carrier \(\Omega_N\), let \(m=2N/5=3a+4\), and sample uniformly
among its \(k\)-subsets. Then
\[
 \Cov(\nu_3,\nu_q)<0\quad(2\le k\le2a+4),
\]
and the covariance is zero for the other admissible cardinalities.
\end{proposition}
\begin{proof}
Each class modulo the odd modulus \(q\) has a unique odd representative
\(h\in[1,2q)\). Its odd terms below \(N\) are
\(h,h+2q,\ldots\); remove the terms divisible by \(5\).
Since \(N>32q\), each such progression has at least sixteen terms.

Any fifteen consecutive terms traverse all odd residues modulo \(30\).
Deleting the multiples of \(5\) leaves four terms in each row modulo \(3\).
Thus adding a full period adds \((4,4,4)\), and \(b=\min v\ge4\).
Before deletion, the three row counts differ by at most one.
Deleted indices have step five; their row counts also differ by at
most one. Hence the surviving row counts have spread at most two.

It remains to check the oriented exception in \eqref{eq:criterion}.
Remove all full periods, leaving a residual progression of length
\(0\le L\le14\). For its normalized profile to be \((0,0,2)\),
row \(2\) must have one extra term before deletion and one fewer
deleted term than each other row. Consequently
\(L\in\{1,4,7,10,13\}\), the progression starts in row \(2\), and
there are exactly two deleted terms, neither in row \(2\).
Indeed the pre-deletion counts can differ by only one, the
deletion counts by only one, and at most three deletions are
possible in fourteen terms.

Write the deleted indices as \(c,c+5,\ldots\), \(0\le c\le4\).
Lengths \(1,4\) allow at most one deletion.
For length \(7\), two deletions require \(c=0\) or \(1\), forcing
a deleted row-\(2\) index at \(0\) or \(6\).
For length \(13\), exactly two require \(c=3\) or \(4\), forcing
a deleted row-\(2\) index at \(3\) or \(9\).
At length \(10\), only \(c=2\) avoids row \(2\), leaving \((2,2,4)\).
Restoring full periods gives
\[
 v=(2+4z,2+4z,4+4z).
\]
The original length is at least sixteen, so \(z\ge1\) and \(b\ge6\).

Finally,
\[
 |v|\le\left\lceil\frac{N}{2q}\right\rceil
 \le\frac{15a+48}{28}\le a\qquad(a\ge4),
\]
where the middle inequality uses \(q\ge7\), \(N=(15a+20)/2\), and
\(\lceil x\rceil\le x+1\). For \(a\ge40\), Theorem~\ref{thm:criterion}
applies.

If \(a<40\), the constraints \(N\equiv10\pmod{30}\),
\(224<N<310\), and \(N>32q\) leave only
\((N,q)=(250,7),(280,7)\).
Their \(101\) and \(113\) possible cardinalities are checked exactly
by \eqref{eq:empty}; the fractions are in
\texttt{small\_cutoffs.json}. They are strictly negative in the
asserted range and zero outside it.
\end{proof}

For \(N=10^d\), these finite extended exceptions are unnecessary:
at \(d=2\) there is no admissible \(7\le q<N/32\), while at
\(d\ge3\) one has \(a\ge132\).

\section{The matching window and exceptional families}\label{sec:matching}
Suppose \(N/4\le q<N/2\), and put \(t=N-2q\).
Every nonsingleton column is a pair with difference \(2q\),
since three odd terms would span at least \(4q\ge N\).
The pairs are exactly
\begin{equation}\label{eq:pairs}
 \{h,h+2q\},\qquad 1\le h<t,\qquad \gcd(h(t-h),10)=1.
\end{equation}
Let \(M\) be their number and let \(\gamma\) count their endpoints
in row \(2\). No pair meets a row twice, because \(3\nmid q\).
Reflection \(h\mapsto N-h\) exchanges rows \(0,1\) and fixes row \(2\),
so the numbers of pairs of each type are
\[
 e_{01}=M-\gamma,\qquad e_{02}=e_{12}=\gamma/2.
\]
The pair centered at \(N/2\), with endpoints \(N/2-q,N/2+q\),
has both endpoints coprime to \(10\); hence \(M\ge1\).

Put
\begin{equation}\label{eq:matchingparams}
 b=2a+2,\quad T=m(m-1),\quad B=b(b-1),\quad S=s(s-1),
\end{equation}
and
\begin{equation}\label{eq:matchingpoly}
 L=T\gamma-2MB,\quad V=M(T-S)-T\gamma,\quad
 P(k)=L(s-k)(s-k-1)+BV.
\end{equation}

\begin{lemma}\label{lem:matching}
For \(2\le k\le s\),
\begin{equation}\label{eq:matching}
 C(d,q,k)=\frac{k(k-1)\binom sk}{\binom mk\,STB}\,P(k).
\end{equation}
Moreover
\[
 P(2)=-MB(3a^2+5a+4)<0,\qquad P(s-1)=P(s)=BV.
\]
\end{lemma}
\begin{proof}
Let \(E_r\) indicate an empty row and \(I_e\) the event that both
endpoints of pair \(e\) are selected.
Then \(\nu_3=3-\sum_r E_r\) and \(\nu_q=k-\sum_e I_e\).
Counting pairs wholly outside each row yields
\[
 C=
 \frac{\gamma\binom{b-2}{k-2}+(M-\gamma)\binom{s-2}{k-2}}{\binom mk}
 -\frac{Mk(k-1)}{T\binom mk}
       \left(2\binom bk+\binom sk\right).
\]
Use
\[
 \binom{n-2}{k-2}=\binom nk\frac{k(k-1)}{n(n-1)},\qquad
 \binom bk=\binom sk\frac{(s-k)(s-k-1)}S
\]
to obtain \eqref{eq:matching}. Substitution at \(k=2,s-1,s\)
gives the endpoint formulas.
\end{proof}

The restriction \(\gcd(q,30)=1\) makes \(t=r+60j\), \(j\ge0\),
with base counts as in Table~\ref{tab:matching}.

\begin{table}[htbp]
\centering
\begin{tabular}{@{}rrrr@{}}
\toprule
\(r\)&\(M_r\)&\(\gamma_r\)&\(D_r=2M_r-3\gamma_r\)\\
\midrule
2&1&0&2\\
6&1&0&2\\
14&5&2&4\\
18&5&4&-2\\
26&7&4&2\\
38&11&6&4\\
42&13&8&2\\
54&17&12&-2\\
\bottomrule
\end{tabular}
\caption{The periodic pair count. For \(t=r+60j\),
\(M=M_r+18j\) and \(\gamma=\gamma_r+12j\).}\label{tab:matching}
\end{table}

The base counts follow directly from \eqref{eq:pairs}.
Over each interval of sixty possible starts there are eighteen
admissible starts, six in each residue class modulo \(3\):
admissibility allows three residues modulo \(10\).
Since the two endpoints occupy distinct rows, these starts give
twelve additional endpoints in row \(2\).
This proves the period formula, without extrapolation from samples.

Since \(X=(s-k)(s-k-1)\) decreases to zero, \(P\) is affine in \(X\).
If \(V<0\), both endpoint values of \(P\) are negative, so all its
values in the range are negative. To have \(V\ge0\), one must have
\begin{equation}\label{eq:ratio}
 \frac{\gamma}{M}\le\frac{T-S}{T}
 =\frac{5a^2+7a}{9a^2+21a+12}<\frac59.
\end{equation}
Consequently \(D_r=2M-3\gamma>M/3\), or \(M<3D_r\).
Table~\ref{tab:matching} leaves only \(t=2,6,14,38\).
For \(t=2,6\), \(V>0\); for \(t=14\),
\[
 V=7a^2-7a-24>0\qquad(a\ge12);
\]
for \(t=38\),
\[
 V=a^2-49a-72,
\]
which is negative at \(a=12\) and positive at every \(a\ge132\).
Thus \(V=0\) never occurs at a decimal depth.

Substituting these four pair counts into \eqref{eq:matchingpoly}
and dividing by positive factors gives exactly the three sign
expressions in Theorem~\ref{thm:decimal}. For \(t=2,6\) the factor is
\(B\); for \(t=14,38\) it is \(2(a+1)\).
Their coefficients of \(X\) are negative.
The sequence \(X\) is strictly decreasing up to \(k=s-1\),
and equals zero at both \(s-1,s\), where the sign is positive.
This proves the single transition and the bound of one integer zero.

\section{Intermediate moduli by affine bands}\label{sec:bands}
We now treat \(N/32\le q<N/4\).
For \(d\ge3\), the proof reduces to fourteen finite certificate families.

\subsection{Column multiplicities}
Put
\[
 \ell=\left\lceil\frac{N}{2q}\right\rceil,\qquad
 \frac{N}{2\ell}\le q<\frac{N}{2(\ell-1)},\qquad 3\le\ell\le16.
\]
For the odd representative \(h\in[1,2q)\) of a class, the odd
progression has length \(\ell\) when
\(h<t=N-2(\ell-1)q\), and length \(\ell-1\) otherwise.
Fix \(u=q\bmod30\) among the eight units modulo \(30\), and put
\[
 t_0=(10-2(\ell-1)u)\bmod30,\qquad b_0=2u\bmod30.
\]
For each odd \(r\in\{1,3,\ldots,29\}\), the numbers of representatives
in the two length ranges are
\begin{align}
 w_{\ell,r}
   &=\frac{N-2(\ell-1)q-t_0}{30}+\one_{\{r<t_0\}},\label{eq:wlong}\\
 w_{\ell-1,r}
   &=\frac{-N+2\ell q+t_0-b_0}{30}
             +\one_{\{r<b_0\}}-\one_{\{r<t_0\}}.\label{eq:wshort}
\end{align}
The column profile from a length \(L\) and a start \(r\) is
\begin{equation}\label{eq:bandvector}
 v_i^{(L,u,r)}
 =|\{0\le j<L:r+2ju\equiv i\pmod3,\ 5\nmid r+2ju\}|.
\end{equation}
There are thirty length/start possibilities.
Aggregating their weights gives
\begin{equation}\label{eq:affine}
 n_v=\frac{A_{\ell,u,v}N+B_{\ell,u,v}q+C_{\ell,u,v}}{30}.
\end{equation}
The checker reconstructs these integer triples from
\eqref{eq:wlong}--\eqref{eq:bandvector}. No precomputed column
multiplicity is trusted without comparison.

\subsection{Sign certificates}
For a fixed band take \(J=\max_v|v|\) and
\(D=\max_{v,r}\delta_r(v)\), over all its residue cases.
The numerator of \eqref{eq:kernelcov}, multiplied by \(30\), is
\begin{equation}\label{eq:bandraw}
 P_{\ell,u}(a,q,k)=
 \sum_v(A_{\ell,u,v}N+B_{\ell,u,v}q+C_{\ell,u,v})K_v(a,k),
 \qquad N=(15a+20)/2.
\end{equation}
It is affine in \(q\).
At each formal endpoint \(q=N/(2\ell)\) and
\(q=N/(2(\ell-1))\), the certificate establishes
\begin{equation}\label{eq:bandcert}
 P_{\ell,u}=k(k-1)R_{\ell,u},\qquad
 R_{\ell,u}\bigm|_{k=2+(2a+2)x}
 =\sum_{i=0}^{J} B_i(a)\binom Ji x^i(1-x)^{J-i}.
\end{equation}
For a positive integer scale \(L\), it records every coefficient of
\(-L B_i(132+y)\), of degree at most \(J+D-1\).
All \(35\,664\) stored integers are strictly positive.
It follows that \(B_i(a)<0\) for every \(a\ge132\).
By Lemma~\ref{lem:certificate}, \eqref{eq:bandcert} is strictly
negative at each endpoint for \(2\le k\le s\).
Affinity in \(q\) proves the same strict sign between them.

These endpoints are evaluations of formal polynomials.
They need not satisfy the residue condition \(q\equiv u\pmod{30}\),
and their column weights need not describe an actual partition.
Affinity requires only the two polynomial endpoint values.

The raw total degree of an endpoint polynomial in \((a,k)\)
is at most \(J+D+1\), and its \(k\)-degree is at most \(J+2\),
by Section~\ref{sec:kernel} and the affine multiplicity factor.
After the substitution for \(k\), the degree bounds in \((a,x)\)
are therefore \(J+D+1,J+2\).
The stored right side of \eqref{eq:bandcert}, including \(k(k-1)\),
has these same separate bounds. Equality at
\(J+D+2\) consecutive \(a\)-values starting at \(132\) and
\(J+3\) nodes \(x=j/(J+2)\) proves the full identity.
In particular, the factor \(k(k-1)\) and the Bernstein
representation are verified, not merely inferred from vanishing
at a few sample sizes.

Table~\ref{tab:bands} lists all dimensions.
Since \(J\le13\), \(D\le11\), and \(a\ge132\), the denominator
in \eqref{eq:kernelcov} is positive.
Thus the verified negative numerator proves negative covariance.

\begin{table}[htbp]
\centering
\begin{tabular}{@{}rrrrr@{}}
\toprule
\(\ell\)&\(J\)&\(D\)&Stored integers&Grid equalities\\
\midrule
3&3&4&448&864\\
4&4&5&720&1232\\
5&4&6&800&1344\\
6&5&6&1056&1664\\
7&6&7&1456&2160\\
8&7&7&1792&2560\\
9&8&8&2304&3168\\
10&8&8&2304&3168\\
11&9&9&2880&3840\\
12&10&9&3344&4368\\
13&11&10&4032&5152\\
14&12&10&4576&5760\\
15&12&10&4576&5760\\
16&13&11&5376&6656\\
\midrule
Total&&&35664&47696\\
\bottomrule
\end{tabular}
\caption{Complete band certificates. Grid counts include all eight
unit residues and both formal endpoints.}\label{tab:bands}
\end{table}


\subsection{The finite depth and the final join}
At \(d=2\), the moduli \(7\le q<N/4\) coprime to \(30\) are
\(7,11,13,17,19,23\).
For these six moduli, the \(246\) fractions from \eqref{eq:empty},
covering all \(41\) cardinalities, are in \texttt{depth2.json}.
The checker recomputes each fraction and compares it with the stored
value. They are strictly negative for \(2\le k\le28\) and zero otherwise.

\begin{proof}[Completion of Theorem~\ref{thm:decimal}]
Section~\ref{sec:empty} proves every deterministic zero.
Section~\ref{sec:matching} proves the classification on
\(N/4\le q<N/2\).
For \(d\ge3\), Section~\ref{sec:bands} proves negativity on
\(N/32\le q<N/4\), while Proposition~\ref{prop:small} proves it for
\(7\le q<N/32\).
The first interval includes its lower endpoint and the second
uses a strict upper endpoint, so the join has no gap.
The exhaustive depth-two calculation completes the remaining cases.
\end{proof}

\section{The arithmetic meaning: two local sieve factors}\label{sec:sieve}
For fixed \(H\subseteq\Omega_N\), consider the translated block forms
\[
 Nb+h,\qquad h\in H,\qquad b\in\mathbb Z.
\]
If \(p\) is prime and \(p\nmid10\), multiplication by \(N\) is
invertible modulo \(p\).
The set of block indices \(b\bmod p\) for which at least one form
is divisible by \(p\) therefore has exactly \(\nu_p(H)\) elements.
Consequently
\begin{equation}\label{eq:localfactor}
 L_p(H)=1-\frac{\nu_p(H)}p
\end{equation}
is the exact fraction of block indices avoiding this local obstruction.
This assertion is a finite residue count, with no distributional
assumption about primes.

\begin{corollary}\label{cor:sieve}
For the uniform \(k\)-subset model and a prime \(q\ge7\),
\begin{equation}\label{eq:sievecov}
 \Cov_H(L_3(H),L_q(H))=\frac{C(d,q,k)}{3q}.
\end{equation}
Equivalently,
\[
 \E_H[L_3L_q]=\E_H[L_3]\E_H[L_q]+\frac{C(d,q,k)}{3q}.
\]
Thus Theorem~\ref{thm:decimal} gives the sign of the exact correction
to the product of these two mean local factors.
\end{corollary}
\begin{proof}
Both local factors are affine functions of their occupancies.
Their covariance is the product of the two slopes times
\(\Cov(\nu_3,\nu_q)\).
\end{proof}

For fixed \(k\), division of each local factor by
\((1-1/p)^k\) is a positive deterministic normalization, so it
preserves the sign in \eqref{eq:sievecov}. Such normalized factors
are the ones appearing in the usual singular-series framework
\cite{kuperberg}. For the block forms here, the primes \(2,5\)
divide no form and give constant local root counts zero.
One must not replace this block model by unrestricted translates
\(b+h\): their local factors at \(2,5\) are different.

For example, at \(N=100\), every prime \(p>50\) acts injectively
on the carrier, so \(\nu_p(H)=k\). Shape dependence of the block
local factors is confined to the thirteen primes
\[
 3,7,11,13,17,19,23,29,31,37,41,43,47.
\]
This finiteness does not turn the two-factor covariance theorem
into a formula for their entire product.

There is also a concrete obstruction to substituting actual primes
for uniform subsets. A uniform subset can occupy all three residue
classes modulo \(3\), giving \(L_3=0\).
A set of offsets whose forms are all primes greater than \(3\)
in one fixed block must avoid the forbidden class modulo \(3\),
so it has \(\nu_3\le2\).
The sampling law in this article is not the law of those prime sets.
No prime prediction, prime-pattern asymptotic, or bound on prime gaps
is deduced.

\section{Consequences, examples and further questions}\label{sec:scope}

\subsection{Strict transitions for two exceptional gaps}

\begin{corollary}[No nontrivial zeros for gaps \(7\) and \(19\)]
\label{cor:nonzero}
For \(2\le k\le s\), the covariance in Theorem~\ref{thm:decimal}
does not vanish when \(q=N/2-7\), \(d\ge2\), or when
\(q=N/2-19\), \(d\ge3\).
In each family its transition from negative to positive is strict
at integer sample sizes. The deterministic zero ranges are unchanged.
\end{corollary}
\begin{proof}
Write \(f_7(a)=(2a+1)(7a^2-7a-24)\) and \(D_7=11a-2\).
Since \(X=(s-k)(s-k-1)\) is integral, a zero would require
\(D_7\mid f_7(a)\). The polynomial identity
\[
 11^3f_7(a)
 =(11a-2)(1694a^2-539a-6753)-45450
\]
then forces \(D_7\mid45450\).
For \(d\ge5\), one has \(D_7\ge146650>45450\).
At \(d=2,3,4\), its values are \(130,1450,14650\), and the
remainders of \(45450\) are respectively \(80,500,1500\).
Thus no admitted depth satisfies the necessary divisibility.

For gap \(19\), put \(f_{19}(a)=(2a+1)(a^2-49a-72)\) and
\(D_{19}=17a-14\). Now
\[
 17^3f_{19}(a)
 =(17a-14)(578a^2-27557a-78471)-1452330,
\]
so a zero requires \(D_{19}\mid1452330\).
For \(d\ge6\), \(D_{19}\ge2266630>1452330\).
At \(d=3,4,5\), the divisors are \(2230,22630,226630\), with
remainders \(600,4010,92550\). None divides the constant.
The endpoint signs and single transition follow from
Theorem~\ref{thm:decimal}.
\end{proof}

For gaps \(1\) and \(3\), put \(j=s-k\).
The common zero equation \(2j(j-1)=a(5a+7)\) has an admissible
integer solution exactly when
\begin{equation}\label{eq:remainingzeros}
 D_d=1+10a^2+14a
\end{equation}
is a square. If so, its square root is odd and
\(j=(1+\sqrt{D_d})/2\); the endpoint signs in
Theorem~\ref{thm:decimal} place this root in the required interval.
For \(n=d-1\ge1\), the same condition is equivalent to the existence of
an integer \(y\) satisfying
\[
 y^2=160\,10^{2n}-152\,10^n+1.
\]
Indeed, the right-hand side equals \(9D_d\), so every such \(y\)
is divisible by \(3\). We leave this exponential Diophantine
question open.

\subsection{Overlapping admissible test sets}

\begin{corollary}[Nonnegative weighted occupancy]\label{cor:weighted}
Let \(E=R_0\sqcup R_1\sqcup R_2\), where
\((|R_0|,|R_1|,|R_2|)=(a+2,a+2,a)=\rho\), for an integer \(a\ge40\).
For a uniform \(k\)-subset \(H\subseteq E\), let \(U(H)\) count
the rows meeting \(H\).
Let \(B_1,\ldots,B_t\subseteq E\) be a finite nonempty family
whose oriented intersection vectors
\[
 v_i=(|B_i\cap R_0|,|B_i\cap R_1|,|B_i\cap R_2|)
     =b_i\one+e_i,\qquad b_i=\min_r(v_i)_r,
\]
each satisfy \eqref{eq:criterion}.
For fixed real weights \(w_i\ge0\) with \(\sum_iw_i>0\), one has
\[
 \Cov\left(U(H),\sum_{i=1}^t
 w_i\one_{\{H\cap B_i\ne\varnothing\}}\right)<0
 \qquad(2\le k\le2a+4).
\]
The covariance is zero at every other allowed cardinality
\(0\le k\le3a+4\). The sets \(B_i\) may overlap and need not cover \(E\).
\end{corollary}
\begin{proof}
Fix an admissible \(B\) with intersection vector \(v\), and let
\(m=3a+4\), \(j=|B|=|v|\).
For \(2\le k\le s=2a+4\), the avoidance identity
\eqref{eq:conditional} gives
\[
 \Cov(U(H),\one_{\{H\cap B\ne\varnothing\}})
 =\frac{\binom{m-j}{k}}{\binom mk}
   \bigl(F(\rho-v;k)-F(\rho;k)\bigr).
\]
Its derivation uses only one set \(B\), not a column partition.
The argument in the proof of Theorem~\ref{thm:criterion},
before summing over columns, establishes
\(F(\rho-v;k)<F(\rho;k)\) for each vector satisfying
\eqref{eq:criterion}. In particular, the orientation of \((0,0,2)\)
and its extra bound on \(b\) are retained.
The prefactor is positive since \(m-j\ge2a+4\ge k\).
Every positively weighted term therefore has negative covariance;
linearity gives the assertion. For \(k=0,1\), or \(k>2a+4\),
the row count \(U\) is deterministic.
\end{proof}

Taking the \(B_i\) to be a column partition and all weights equal to
one recovers Theorem~\ref{thm:criterion}. Alternatively, weights may
be assigned only to admissible columns of a larger partition.
No signed-weight, stochastic negative-association, or nonuniform
sampling statement is asserted.

\subsection{Examples and remaining questions}
The sign changes are already visible in the smallest decimal carrier:
\[
 C(2,47,8)=\frac{1771}{20974005}>0.
\]
At a common cardinality \(k=15\), one has
\[
 C(2,41,15)=-\frac{51037}{134914728}<0,\qquad
 C(2,43,15)=\frac{161}{10939032}>0.
\]
The row sizes are the same, and both column partitions have five
pairs and thirty singletons, hence the same marginal law for
\(\nu_q\). Their intersections with the three rows differ.
Thus marginal occupancy laws alone cannot determine the covariance.

The matroid interpretation also has a limit.
For \(N=100,q=7\), the two rank functions are not a matroid
perspective in either direction. Adding \(17\) to \(\{3\}\)
increases the modulo-\(3\) rank by one and the modulo-\(7\) rank
by zero. Adding \(7\) to \(\{1\}\) does the reverse.
These two examples violate the respective rank-increment inequalities
for a perspective \cite[p.~304]{lasvergnas}.
The sign theorem is therefore not presented as an application of
a perspective-specific positivity statement.

Several questions remain distinct from the proved results.
The thresholds \(h\ge16\) and \(a\ge40\) in the urn inequalities
may admit improvement, but sharpness is not needed here.
An elementary proof eliminating their coefficient certificates
would shorten the computational dependency.
After Corollary~\ref{cor:nonzero}, the integer-zero question is
restricted to gaps \(1\) and \(3\), equivalently the square condition
\eqref{eq:remainingzeros}.
Other bases and other pairs of moduli require separate analysis.
Finally, a useful transfer to prime patterns would need a
controlled full-product average or a justified nonuniform law on
offset sets. Pairwise covariance supplies only one part of that problem.

\section{Certificate specification and reproduction}\label{sec:certificates}
The companion folder is self-contained and uses Python 3.10 or later.
It contains eighteen JSON certificate files:
fourteen bands, two urn files, and two finite-boundary files.
A manifest pins all twenty-six payload files by byte count and SHA-256;
the manifest itself is the twenty-seventh file.
The original certificate bytes are retained, with their source-member
identities in \texttt{PROVENANCE.json}.
The checker requires no original ZIP, symbolic generator, historical
prime dataset, or internal project document.

\subsection{Encodings and independent reconstruction}
For a band, the certificate declares \(\ell,J,D\), a positive scale,
the shift \(132\), affine type tables, endpoint denominators and every
Bernstein index.
The checker derives \(J,D\), the type alphabet and the affine
multiplicities independently from \eqref{eq:bandvector}.
It requires exactly all eight residues and both endpoints.
The integer array for index \(i\) lists the coefficients of
\(-L B_i(132+y)\) in ascending powers of \(y\).

For the urn and anchor files, the arrays are the coefficients of
\(\sigma L B_i(t_0+y)\), again in ascending powers.
Here \((\sigma,t_0)=(1,16)\) for the urn increment and
\((-1,40)\) for the anchors.
Thus the reconstructed Bernstein coefficient is
\(\sigma/L\) times the stored power polynomial.
Array dimensions, positive scale, domain shift, sign, and positive
constant terms are checked explicitly.
The raw polynomials come from the displayed formulas in the proofs,
not from the symbolic generator that produced the arrays.

The complete identity workload is
\[
 47\,696+890+3166=51\,752
\]
exact rational equalities, using \(38\,784\) stored integer coefficients.
In addition, the finite boundaries require \(246+214=460\) exact
fractions. Direct probability comparisons and adverse examples
provide separate checks of the algebra-to-probability conversion.
They supplement, rather than replace, the global polynomial identities.

\newpage
\subsection{How to run and interpret the checks}
From the companion directory, run
\begin{verbatim}
python -B verify.py
python -B -O verify.py
python -B -m unittest -v test_companion
python -B -O -m unittest -v test_companion
\end{verbatim}
The verifier writes nothing and prints a deterministic JSON receipt.
If output is redirected, the destination should be outside the companion
directory, whose inventory is checked exactly.
The normal and optimized receipts must agree byte for byte.
The suite contains 17 unit tests: 16 negative controls and one positive
inventory check.
The test suite checks rejection of omitted cases, corrupted coefficients,
altered geometry, weakened degree declarations, wrong domains and
missing files; it does not rely on Python assertions being enabled.
A clean copy of the folder suffices to repeat every check.

The certificate map identifies each proof dependency and grid size.
Hashes establish consistency with that particular package, not a
trusted digital signature. A successful checker run does not prove
that an arbitrary revised manuscript correctly invokes the identities:
the written domain, positivity and reduction arguments remain part
of the mathematical proof.

\section*{Preparation note}
AI assistance contributed to mathematical exploration, proof development,
code preparation, verification and drafting. The author supplied research
packages produced with an external AI agent; its exact model and version
are not documented in the available record. Integration, manuscript
preparation and computational checks used the Codex assistant.
A separate AI reading examined an earlier version. This is not human
peer review, and complete human verification is not claimed here.
The companion records certificate provenance and checker adaptations.
Exact computational replay supports the stated identities; it does not
establish novelty or replace review of the written arguments.

\begin{thebibliography}{9}
\bibitem{benhamou}
A.~Ben-Hamou, Y.~Peres and J.~Salez,
\emph{Weighted sampling without replacement},
arXiv:1603.06556 (2016).
\url{https://arxiv.org/abs/1603.06556}.

\bibitem{joagdev}
K.~Joag-Dev and F.~Proschan,
Negative association of random variables, with applications,
\emph{Annals of Statistics} \textbf{11} (1983), no.~1, 286--295.
\href{https://doi.org/10.1214/aos/1176346079}{doi:10.1214/aos/1176346079}.

\bibitem{klm}
A.~Klenke and L.~Mattner,
Stochastic ordering of classical discrete distributions,
\emph{Advances in Applied Probability} \textbf{42} (2010), no.~2, 392--410.
\href{https://doi.org/10.1239/aap/1275055235}{doi:10.1239/aap/1275055235}.
Preprint:
\url{https://arxiv.org/abs/0903.1361v2}.

\bibitem{kuperberg}
V.~Kuperberg,
Sums of singular series along arithmetic progressions and with
smooth weights,
\emph{International Journal of Number Theory} \textbf{21} (2025), no.~1, 53--74.
\href{https://doi.org/10.1142/S1793042125500046}{doi:10.1142/S1793042125500046}.
Preprint:
\url{https://arxiv.org/abs/2301.06095}.

\bibitem{lasvergnas}
M.~Las Vergnas,
The Tutte polynomial of a morphism of matroids.
4. Computational complexity,
\emph{Portugaliae Mathematica} \textbf{64} (2007), no.~3, 303--309.
\url{https://ems.press/content/serial-article-files/44572?nt=1}.
\end{thebibliography}
\end{document}
